%!TEX TS-program = Arara
% arara: pdflatex: { draft : yes }
% arara: nomencl
% arara: biber
% arara: pdflatex

%-------------------------------------------------------------------------------
%Dokumentenklasse und Schriftarten
%-------------------------------------------------------------------------------
\documentclass[12pt, a4paper, ngerman,toc=chapterentrydotfill, listof=totoc,bibliography=totoc,captions=figureheading,hidelinks,listof=nochaptergap,fleqn]{scrreprt} 
\usepackage[T1]{fontenc} 
\usepackage[utf8]{inputenc}
\usepackage{palatino}%schönere schrift 
\usepackage{microtype} % Mikrotypografie
\usepackage[ngerman]{babel}%Sprache und Silbentrennung
\usepackage{paralist}

\usepackage[onehalfspacing]{setspace}%Zeilenabstand auf 1,5 stellen
\usepackage[a4paper, left=4cm, right=2cm, top=1.5cm, bottom=2cm, headsep=2cm, includehead=true, footskip=0cm, includefoot=true, headheight = 0.5cm]{geometry} 


%-------------------------------------------------------------------------------
%Mathematische Ausdrücke
%-------------------------------------------------------------------------------
%\usepackage{amsmath}%Mathematische Ausdrücke
\counterwithout{equation}{chapter}%Nummerierung aufsteigend ohne Kapitelnummer

% Mathe Einrückung
\setlength{\mathindent}{1cm}

%-------------------------------------------------------------------------------
%Link Umbrueche im Verzeichnis
%-------------------------------------------------------------------------------
\usepackage{hyperref}%zum Setzen von PDF Daten und umbrechen von Links
\sloppy

%-------------------------------------------------------------------------------
%Meta Infos laden
%-------------------------------------------------------------------------------
%-----------------------------------
% Meta Informationen zur Arbeit
%-----------------------------------

% Autor
\newcommand{\myAutor}{Karl Toffel}
% Adresse
\newcommand{\myAdresse}{Musterstrasse 3 \\ \> \> 12345 Musterstadt}
% Titel der Arbeit
\newcommand{\myTitel}{\LaTeX-Vorlage - FOM Leitfaden}
% Betreuer
\newcommand{\myBetreuer}{Prof. Dr. Max Mustermann}
% Lehrveranstaltung
\newcommand{\myLehrveranstaltung}{Mustermodul}
% Matrikelnummer
\newcommand{\myMatrikelNr}{123456}
% Ort
\newcommand{\myOrt}{Köln}
% Datum der Abgabe
\newcommand{\myAbgabeDatum}{30. Januar 2019}
% Semesterzahl
\newcommand{\mySemesterZahl}{3}
% Name der Hochschule
\newcommand{\myHochschulName}{FOM Hochschule für Oekonomie \& Management}
% Standort der Hochschule
\newcommand{\myHochschulStandort}{Standort Köln}
% Studiengang
\newcommand{\myStudiengang}{Wirtschaftsinformatik}
% Art der Arbeit
\newcommand{\myThesisArt}{Seminararbeit}
% Zu erlangender akademische Grad
\newcommand{\myAkademischerGrad}{Bachelor of Science (B. Sc.)}
% Firma
\newcommand{\myFirma}{Mustermann GmbH}

%---------------------------------------
% PDF Meta setzen
%---------------------------------------
\hypersetup{
    pdfinfo={
        Title={\myTitel},
        Subject={\myStudiengang},
        Author={\myAutor},
        Build=1.1
    }
}

% Mathe 

\usepackage{amsmath}
\usepackage{amssymb}

%-------------------------------------------------------------------------------
%Abbildungen einbinden
%-------------------------------------------------------------------------------
\usepackage{graphicx}
\usepackage{float}%Festlegen von Float Objects mit [H]
\usepackage{chngcntr}%Paket für Abbildungsnummerierung ändern
\counterwithout{figure}{chapter}%ändern der Nummerierung zu aufsteigend
\counterwithout{table}{chapter}%ändern der Nummerierung zu aufsteigend
\usepackage[labelfont=bf,position=top, belowskip = 6pt, aboveskip=12pt]{caption}%Abstand zum Text festlegen

%-------------------------------------------------------------------------------
%Zitate und Literaturverzeichnis
%-------------------------------------------------------------------------------
\usepackage[babel,german=guillemets]{csquotes} %Festlegen der Anführungszeichen
\usepackage[backend=biber, style=ext-authoryear, isbn=false, maxcitenames=2, maxbibnames=999, mergedate=false, innamebeforetitle, dashed=false, autocite=footnote, mincrossrefs=1, autopunct=false]{biblatex}%Bibliotheksverwaltung (siehe goLatex FOM)

\usepackage{blindtext}%Paket für Blindtext Generierung

%-------------------------------------------------------------------------------
%Abkuerzungsverzeichnis
%-------------------------------------------------------------------------------
\usepackage[intoc]{nomencl}
\renewcommand{\nomname}{Abkürzungsverzeichnis}%umbenennen
\setlength{\nomlabelwidth}{.15\textwidth}%Breite des Tabstops festlegen
%\renewcommand{\nomlabel}[1]{#1 \dotfill}%Füllen mit Punkt
\setlength{\nomitemsep}{-\parsep}
\makenomenclature

%-------------------------------------------------------------------------------
%Weitere Einstellungen
%-------------------------------------------------------------------------------
\setlength{\parindent}{0em}%Einrückung nach Absatz auf 0 setzen}%Seitenränder gesetzt
\setlength{\parskip}{6pt}%Absatzabstand
%-----------Kopf/Fußzeile--------------

\usepackage[automark, headsepline=.4pt]{scrlayer-scrpage}%Paket für Kopfzeile
\newpairofpagestyles{myheading}%Kopfzeile definieren
{
	\ihead{\leftmark}%Kapitelname links in Kopfzeile
	\chead{}%nichts in der Mitte
	\ohead{\thepage}%Seitenzahl rechts in Kopfzeile
}

\usepackage[automark, headsepline=.4pt]{scrlayer-scrpage}%Paket für Kopfzeile
\newpairofpagestyles{leer}%Kopfzeile definieren
{
	\ihead{}%Kapitelname links in Kopfzeile
	\chead{}%nichts in der Mitte
	\ohead{}%Seitenzahl rechts in Kopfzeile
}


\pagestyle{myheading}%Kopfzeile aktiv machen
\renewcommand*\chapterpagestyle{myheading}%Kopfzeile auch auf Seiten mit neuem Kapitel anzeigen

%----------Kapitelüberschriften formatieren-------

\setkomafont{disposition}{\rmfamily\bfseries}%Alle Überschriften mit Serifen und Fett
\setkomafont{chapter}{\huge}%Kapitelüberschriften in Groß
\setkomafont{section}{\large}%Sektionsüberschriften in Groß

\RedeclareSectionCommand[beforeskip=0pt, afterskip=6pt]{chapter}%12pt vor und 6pt nach Kapitelüberschrift
\RedeclareSectionCommand[beforeskip=12pt, afterskip=6pt]{section}%12pt vor und 6pt nach Kapitelüberschrift

\newcommand{\befehl}[1]{\lstinline{#1}}%Datei mit Zitieranpassungen und Format des Literaturverzeichnisses
\setlength\bibhang{1cm} 

\renewcommand*{\finentrypunct}{} 

\renewcommand*{\newunitpunct}{\addcomma\space} 
\DeclareDelimFormat[bib,biblist]{nametitledelim}{\addcolon\space} 
\DeclareDelimFormat{titleyeardelim}{\newunitpunct} 
\renewcommand*{\mkbibnamefamily}{\mkbibemph} 
\renewcommand*{\mkbibnamegiven}{\mkbibemph} 
\renewcommand*{\mkbibnamesuffix}{\mkbibemph} 
\renewcommand*{\mkbibnameprefix}{\mkbibemph} 
\DeclareDelimAlias{finalnamedelim}{multinamedelim} 

\DeclareNameFormat{family-given:init}{% 
  \usebibmacro{name:family-given} 
    {\namepartfamily} 
    {\namepartgiveni} 
    {\namepartprefix} 
    {\namepartsuffix}% 
  \usebibmacro{name:andothers}} 

\DeclareNameAlias{default}{family-given} 
\DeclareNameAlias{sortname}{default} 
\DeclareNameAlias{labelname}{family-given:init} 

\DeclareFieldFormat{editortype}{\mkbibparens{#1}} 
\DeclareDelimFormat{editortypedelim}{\addspace} 
\DeclareFieldFormat{translatortype}{\mkbibparens{#1}} 
\DeclareDelimFormat{translatortypedelim}{\addspace} 
\DeclareDelimFormat[bib,biblist]{innametitledelim}{\addcomma\space} 

\DeclareFieldFormat*{citetitle}{#1} 
\DeclareFieldFormat*{title}{#1} 
\DeclareFieldFormat*{booktitle}{#1} 
\DeclareFieldFormat*{journaltitle}{#1} 

\newbibmacro*{organization+location}[1]{% 
  \printlist{location}% 
  \setunit*{\locdatedelim}% 
  \printlist{organization}% 
  \newunit} 


\usepackage{xpatch} 
\xpatchbibdriver{online} 
  {\usebibmacro{organization+location+date}} 
  {\usebibmacro{organization+location}} 
  {}{} 

\DeclareFieldFormat[online]{date}{\mkbibparens{#1}} 

\DeclareFieldFormat{url}{<\url{#1}>}% oder {$\langle$\url{#1}$\rangle$} oder \textlangle/\textrangle 
\DeclareFieldFormat{urldate}{\mkbibbrackets{\bibstring{urlseen}\addcolon\space#1}} 
\renewbibmacro*{url+urldate}{% 
  \usebibmacro{url}% 
  \ifentrytype{online} 
    {\setunit{\addspace}% 
     \usebibmacro{date}} 
    {}% 
  \iffieldundef{urlyear} 
    {} 
    {\setunit*{\addspace}% 
     \usebibmacro{urldate}}} 

\renewbibmacro*{date+extradate}{% 
  \printtext[parens]{% 
    \printfield{usera}% 
    \setunit{\printdelim{titleyeardelim}}% 
    \printlabeldateextra}} 

\DefineBibliographyStrings{german}{ 
  nodate    = {{}o.\adddot J\adddot}, 
  andothers = {et\addabbrvspace al\adddot}, 
  urlseen   = {Zugriff}, 
} 

\DeclareSourcemap{ 
  \maps[datatype=bibtex]{ 
    \map{ 
      \step[notfield=translator, final] 
      \step[notfield=editor, final] 
      \step[fieldset=author, fieldvalue={""{o\noexpand\adddot V.}""}]
    } 
    \map{ 
      \pernottype{online} 
      \step[fieldset=location, fieldvalue={o\noexpand\adddot O\noexpand\adddot}] 
    } 
  } 
} 

\renewbibmacro*{cite}{% 
  \iffieldundef{shorthand} 
    {\ifthenelse{\ifnameundef{labelname}\OR\iffieldundef{labelyear}} 
       {\usebibmacro{cite:label}% 
        \setunit{\printdelim{nonametitledelim}}} 
       {\printnames{labelname}% 
        \setunit{\printdelim{nametitledelim}}}% 
     \printfield{usera}% 
     \setunit{\printdelim{titleyeardelim}}% 
     \usebibmacro{cite:labeldate+extradate}} 
    {\usebibmacro{cite:shorthand}}} 

\renewbibmacro*{postnote}{% 
  \setunit{\postnotedelim}% 
  \iffieldundef{postnote} 
    {\printtext{o. S\adddot}} 
    {\printfield{postnote}}}
    
\newcommand{\directAutocite}[3][]{\enquote{#3}\autocite[#1]{#2}}%Command für direkte Zitate%Datei mit Zitieranpassungen und Format des Literaturverzeichnisses
% sauber formatierter Quelltext
\usepackage{listings}
\lstset{numbers=left,
	numberstyle=\tiny,
	numbersep=5pt,
	breaklines=true,
	showstringspaces=false,
	frame=l ,
	xleftmargin=5pt,
	xrightmargin=5pt,
	basicstyle=\ttfamily,%\scriptsize, 
	stepnumber=1,
	keywordstyle=\color{blue},          % keyword style
  	commentstyle=\color{dkgreen},       % comment style
  	stringstyle=\color{mauve}         % string literal style
}

\let\befehl\lstinline%Sauber formatierter Quelltext
\addbibresource{lit.bib}%include der Literatur Datenbank
%\nocite{*}%UNBEDINGT AUSKOMMENTIEREN. DURCH DIESEN BEFEHL WERDEN AUCH NICHT ZITIERTE QUELLEN VERZEICHNET
%\addbibresource{biblatex-examples.bib}%Beispielliteratur einfügen



%-------------------------------------------------------------------------------
%Dokumentbeginn
%-------------------------------------------------------------------------------
\begin{document}%Beginn des Dokumentes
\newcounter{savepage}
\pagenumbering{Roman}%Die ersten Seiten römisch nummerieren

%-------------------------------------------------------------------------------
%Titelblatt
%-------------------------------------------------------------------------------
\begin{titlepage}
	\newgeometry{left=4cm, right=2cm, top=4cm, bottom=2cm}
	\begin{center}
		\textbf{\myHochschulName}\\
		\textbf{\myHochschulStandort}\\
		\vspace{1.5cm}
			\includegraphics[width=3cm]{Abbildungen/fom_logo.png} \\
		\vspace{1cm}
		Berufsbegleitender Studiengang\\
		\myStudiengang, \mySemesterZahl. Semester\\
		\vspace{1cm}
		\textbf{\myThesisArt}\\
		%\textbf{zur Erlangung des Grades eines}\\
		%\textbf{\myAkademischerGrad}\\
		\textbf{im Rahmen der Lehrveranstaltung}\\
		\textbf{\myLehrveranstaltung}\\
		\vspace{1cm}
		über das Thema\\
		\Huge{\myTitel}\\
		\vspace{0.2cm}
	\end{center}
	\normalsize
	\vfill
	\begin{tabbing}
		Links \= Mitte \= Rechts\kill
		Betreuer: \> \> \myBetreuer\\
		\> \> \\

		Autor: \> \> \myAutor\\
		\> \>  Matrikelnr.: \myMatrikelNr\\
		\> \> \myAdresse\\
		\> \> \\
		Abgabe: \> \> \myAbgabeDatum
	\end{tabbing}
\end{titlepage}

%-------------------------------------------------------------------------------
%Verzeichnisse
%-------------------------------------------------------------------------------
\setcounter{page}{2}%TOC als erste Seite
\tableofcontents%Inhaltsverzeichnis einfügen
\listoffigures%Abbildungsverzeichnis einfügen
\listoftables%Tabellenverzeichnis einfügen
\printnomenclature%Abkürzungsverzeichnis einfügen
\cleardoublepage

%-------------------------------------------------------------------------------
%Text
%-------------------------------------------------------------------------------
\setcounter{savepage}{\arabic{page}}%Seitenstand der römischen Nummerierung speichern um später dort weiterzuzählen
\pagenumbering{arabic}%arabische Nummerierung

\input{Textteil/Kapitel/Einleitung}
\chapter{Installation}
\label{ch:installation}
Zur Nutzung von \LaTeX\ müssen zunächst einmal einige Dinge auf dem Rechner installiert werden. Es folgt eine kleine Anleitung, mit der auch das kein Problem sein sollte.

\section{MiKTeX}
MiKTeX ist neben TeX Live und MacTex eine der größten Distributionen, die es von TeX gibt. Für Windows Systeme ist diese sinnvoll und installiert Pakete, welche noch nicht heruntergeladen wurden, automatisch nach.

Der Download Link für MiKTeX ist in den Quellen vorhanden.\autocite[Vgl.][]{miktex} Hierbei kann der gewöhnliche \directAutocite{miktex}{Installer} verwendet werden. Anschließend führt das Installationsprogramm den Nutzer durch den gesamten Installationsprozess. Hierbei muss nichts besonderes beachtet werden.

\section{TexMaker}
TexMaker ist der Editor und damit die Schnittstelle zum Nutzer. Zwar ist es möglich, sämtliche Dokumente in Texteditoren wie Notepad++ zu erstellen und diese über die Eingabeaufforderung zu einem PDF (\nomenclature{PDF}{Portable Document Format}Portable Document Format) umzuwandeln, doch bietet TexMaker angenehme Funktionen für die erleichterte Nutzung.

Auch hier gibt es einen normalen Installer zum Download, welcher alle weiteren Schritte erklärt.\autocite[Vgl.][]{texmaker}

Zum Einrichten von TexMakter muss noch eine weitere Einstellung vorgenommen werden. Da \LaTeX\ zum Bilden von Referenzen, Generieren von Verzeichnissen etc zunächst temporäre Dateien erzeugt, welche erst beim nächsten compilen genutzt werden, müssen eine Reihe von Compile-Befehlen nacheinander aufgerufen werden. Hierzu 
\chapter{Meta-Daten}
Die Meta Daten beschreiben den Studierenden und alles was sonst noch dazu gehört. Um das Projekt einigermaßen übersichtlich zu halten, habe ich diese in eine eigene Datei ausgelagert. Hier werden Commands erzeugt, welche später im Projekt (wie der Titelseite) aufgerufen werden können, um den jeweiligen Bestandteil (wie den Autor) einzufügen. Hier sollten also alle Angaben entsprechend der anzufertigenden Arbeit angepasst werden. Eine Arbeit unter fremdem Namen macht sich wohl nicht so gut bei den meisten Prüfern.
% !TeX root = ../../main.tex

\chapter{Fließelemente}\label{ch:fliesselemente}

Fließelemente sind Elemente, welche sich (wie der Name schon sagt) fließend im Text bewegen. Das bedeutet, dass \LaTeX\ diese dort positioniert, wo es \enquote{denkt}, dass es am besten passt. Dies funktioniert in den meisten Fällen auch recht gut.
Zu Fließelementen gehören vor allen Dingen Abbildungen und Tabellen. 
Um diese der FOM entsprechend einzufügen, muss eine Reihe von Befehlen genutzt werden.

Es gibt einige Felder, welche hier vom Nutzer angepasst werden müssen:

\begin{description}
\item[label]
ist eine Art ID, welche im gesamten Dokument referenziert werden kann. Mehr hierzu in Kapitel \ref{ch:referenzen} auf Seite \pageref{ch:referenzen}.
\item[caption] ist die Überschrift des Fließelements, welche auch in Verzeichnissen angezeigt wird.
\item[width] ist die Breite des Elementes. Sie kann entweder absolut mit z. B. enquote{width = 8cm}, oder relativ mit \enquote{width=0.5\textbackslash linewidth} angegeben werden.
\end{description}

\section{Abbildungen}
Zum Einbinden von Abbildungen werden folgende Befehle innerhalb des Textes genutzt:
\begin{lstlisting}
% Hier wird aufgezeigt, wie man eine Grafik einbindet,auch wenn es hier faelschlicherweise als LaTeX-Befehl angezeigt wird.
\begin{figure}[H] 
   \centering 
   \caption{FOM Logo 2019} 
   \label{fig:fom_logo} 
   \includegraphics[width=0.5\linewidth]{Abbildungen/fom_logo.png} 
   \par\smallskip 
   Quelle: \cite[12]{bsp} 
\end{figure} 
\end{lstlisting}
Dieses Beispiel sieht anschließend folgendermaßen aus:
\begin{figure}[H] 
   \centering 
   \caption{FOM Logo 2019} 
   \label{fig:fom_logo} 
   \includegraphics[width=0.5\linewidth]{Abbildungen/fom_logo.png} 
   \par\smallskip 
   Quelle: \cite[12]{bsp} 
\end{figure} 

\section{Tabellen}
Tabellen werden ähnlich wie Abbildungen eingefügt. 
Hinzu kommt jedoch die Syntax, mit welcher die Tabelle als solche (unabhängig vom Float Element) erzeugt wird. 
Ein Beispiel ist folgendes:

\begin{lstlisting}
\begin{table}[H]
	\centering
	\caption{Tabellenname}
	\label{table:example}
	\begin{tabular}{||c c c c||}
		\hline
		N1 & N2 & N3 & N4
		\\ [0.5ex]
		\hline\hline
		1 & 2 & 3 & 4\\
		5 & 6 & 7 & 8\\
		[1ex]
		\hline
	\end{tabular}
	\par\smallskip 
	Quelle: \cite[13]{bsp}
\end{table}
\end{lstlisting}

Diese Tabelle sieht anschließend folgendermaßen aus:
\begin{table}[H]
	\centering
	\caption{Beispieltabelle}
	\label{table:example}
	\begin{tabular}{||c c c c||}
		\hline
		N1 & N2 & N3 & N4
		\\ [0.5ex]
		\hline\hline
		1 & 2 & 3 & 4\\
		5 & 6 & 7 & 8\\
		[1ex]
		\hline
	\end{tabular}
	\par\smallskip 
	Quelle: \cite[13]{bsp}
\end{table}

In den Parametern des ''tabular''s wird die Anzahl der Spalten und die Darstellung der vertikalen Linien bestimmt. 
Ein Pipe-Symbol steht hierbei für eine vertikale Linie an der entsprechenden Stelle, während ein Buchstabe für eine Spalte steht. Es können folgende Buchstaben genutzt werden:
\begin{description}
	\item[c] steht für eine Spalte mit zentriertem Text.
	\item[l] steht für eine Spalte mit linksbündigem Text.
	\item[r] steht für eine Spalte mit rechtsbündigem Text.
\end{description}

Im Anschluss wird durch \enquote{\textbackslash hline} eine horizonale Linie erzeugt. 
Diese ist in diesem Fall die oberste Linie der Tabelle. 
Die nächste Zeile beschreibt die Spalteninhalte. 
Hier werden in diesem Fall Spaltennamen (N1 etc.) gesetzt, bevor zwei weitere Linien zur Begrenzung der Titelzeile eingefügt werden. 
Darauffolgend werden die restlichen Zellen gesetzt.
% !TeX root = ../../main.tex

\chapter{Referenzen}\label{ch:referenzen}

Referenzen werden genutzt, wenn innerhalb des Textes auf Abbildungen, Tabellen oder Kapitel verwiesen werden soll. 
Um beim Einfügen eines weiteren Kapitels einen Verweis zu verfälschen (da Kapitel X nun etwas anderes ist), wird das Objekt über einen Marker (gesetzt durch \enquote{\textbackslash label\{marker\}}) referenziert. 
Hierzu gibt es zwei wichtige Befehle.

\begin{description}
\item[{\textbackslash nameref\{marker\}}] erzeugt den Namen des Objektes.
\item[{\textbackslash ref\{marker\}}] erzeugt die Objektnummer, also ob es sich um Tabelle 1, 2 oder 3 etc. handelt.
\item[{\textbackslash pageref\{marker\}}] erzeugt die Seitennummer, auf welcher sich das Objekt befindet.
\end{description}

\section{Beispiel}
Hier ist ein kleines Beispiel zur Nutzung von Referenzen.

\begin{lstlisting}
Die Abbildung \nameref{fig:fom_logo} hat die Nummer \ref{fig:fom_logo} und ist auf Seite \pageref{fig:fom_logo}
\end{lstlisting}
Ausgabe: \enquote{Die Abbildung \nameref{fig:fom_logo} hat die Nummer \ref{fig:fom_logo} und ist auf Seite \pageref{fig:fom_logo}}

\section{Warnung}

Dieses Kapitel ist lediglich dazu da, dich daran zu erinnnern, dass es niemals ein Kapitel mit nur einem Unterkapitel geben darf. ;-)
\chapter{Zitieren}
\label{ch:zitieren}
Das Zitieren ist ein elementarer Bestandteil einer jeden wissenschaftlichen Arbeit. Um alle Formvorschriften einzuhalten empfehle ich, lediglich drei verschiedene Befehle zu nutzen:
\begin{description}
	\item[{\textbackslash cite[a][b]\{c\}}] gibt den Fußnoteninhalt nicht in einer Fußnote, sondern direkt im Text aus. Dies wird vor allem bei der Quellenangabe von Abbildungen verlangt.
	\item[{\textbackslash autocite[a][b]\{c\}}] gibt eine gewöhnliche Fußnote mit Quelle aus.
	\item[{\textbackslash directAutocite[a][b]\{c\}\{d\}}] erzeugt ein direktes Zitat mit zugehöriger Fußnote. Die Anführungszeichen werden je nach vorheriger Festlegung gewählt (hier standardmäßig Guillements).
\end{description}
Für die Parameter a, b, c und d gilt:
\begin{description}
	\item[a] ist das Präfix der Fußnote (z. B. \enquote{Vgl.}). Dieser Parameter ist optional.
	\item[b] ist die Seitenzahl, von der die Information(/Zitat) entnommen wurde. Auch dieser Parameter ist optional. Sollte er jedoch weggelassen werden, sollten die Klammern trotzdem geschrieben werden, da ansonsten ein Suffix nicht vom Präfix unterschieden werden kann. In diesem Fall werden die Klammern also lediglich freigelassen.
	\item[c] ist die ID der Quelle.
	\item[d] ist der Text, welcher innerhalb der Anführungszeichen aufgeführt werden soll.
\end{description}

Hier ein Beispiel:
\begin{lstlisting}
\directAutocite[12]{bsp}{Das hier ist ein Zitat}
\end{lstlisting}
Ausgabe: \directAutocite[12]{bsp}{Das hier ist ein Zitat}
% !TeX root = ../../main.tex

\chapter{Kapitel einfügen}

Das Einfügen der Kapitel ist ganz einfach. 
Zunächst muss eine Datei im Verzeichnis Textteil/Kapitel angelegt werden. 
In dieser datei steht der Inhalt des gesamten Kapitels, inklusive \enquote{\textbackslash chapter\{Name des Kapitels\}}. 
Diese wird in der Datei Textteil/textteil.tex per \verb|\input| eingefügt. 
Ein Beispiel hierfür sind die bereits eingefügten Kapitel.
\input{Textteil/Kapitel/Schluss}
% !TeX root = ../../main.tex

%-----------------------------------
% Ehrenwörtliche Erklärung
%-----------------------------------
\chapter*{Ehrenwörtliche Erklärung}
\pagenumbering{gobble} % Keine Seitenzahlen mehr
\renewcommand*\chapterpagestyle{leer}
Hiermit versichere ich, dass die vorliegende Arbeit von mir selbstständig und ohne unerlaubte Hilfe angefertigt worden ist, insbesondere dass ich alle Stellen, die wörtlich oder annähernd wörtlich aus Veröffentlichungen entnommen sind, durch Zitate als solche gekennzeichnet habe. Ich versichere auch, dass die von mir eingereichte schriftliche Version mit der digitalen Version übereinstimmt. Weiterhin erkläre ich, dass die Arbeit in gleicher oder ähnlicher Form noch keiner Prüfungsbehörde/Prüfungsstelle vorgelegen hat. Ich erkläre mich damit einverstanden, dass die Arbeit der Öffentlichkeit zugänglich gemacht wird. Ich erkläre mich damit einverstanden, dass die Digitalversion dieser Arbeit zwecks Plagiatsprüfung auf die Server externer Anbieter hoch geladen werden darf. Die Plagiatsprüfung stellt keine Zurverfügungstellung für die Öffentlichkeit dar.

\par\medskip
\par\medskip

\vspace{5cm}

\begin{table}[H]
	\centering
	\begin{tabular*}{\textwidth}{c @{\extracolsep{\fill}} ccccc}
		\myOrt, \today
		&
		% Hinterlege deine eingescannte Unterschrift im Verzeichnis /abbildungen und nenne sie unterschrift.png
		% Bilder mit transparentem Hintergrund können teils zu Problemen führen
		\includegraphics[width=0.35\textwidth]{Abbildungen/unterschrift.png}\vspace*{-0.35cm}
		\\
		\rule[0.5ex]{12em}{0.55pt} & \rule[0.5ex]{12em}{0.55pt} \\
		(Ort, Datum) & (Eigenhändige Unterschrift)
		\\
	\end{tabular*} \\
\end{table}%Datei mit eigentlichem Text einbinden
\chapter{Anhang}
\renewcommand\thesection{Anhang \arabic{section}:} %https://tex.stackexchange.com/questions/3177/how-to-change-the-numbering-of-part-chapter-section-to-alphabetical-r
\addtocontents{toc}{\protect\setcounter{tocdepth}{-1}}
%-------------------------------------------------------------------------------

\begin{minipage}{\textwidth}
	\section{Versuch Nummer Eins}
	\label{anh:versuch_nummer_eins}
	\includegraphics[width=0.5\linewidth]{Abbildungen/fom_logo.png} 
	\par\smallskip
	Quelle: \cite{logo}
\end{minipage}

\addtocontents{toc}{\protect\setcounter{tocdepth}{1}}
\renewcommand\thesection{\thechapter.\arabic{section}}
%-------------------------------------------------------------------------------
%Verzeichnisse
%-------------------------------------------------------------------------------
\pagenumbering{Roman}%römische Nummerierung
\setcounter{page}{\thesavepage}%Seitencounter auf die gespeicherte Zahl setzen

%--------FÜR HAUSARBEIT NACH LEITFADEN----------------------
\addchap{Literaturverzeichnis}
\printbibliography[nottype=online,heading=none]%Literaturverzeichnis der NICHT-Onlinequellen drucken und ins Inhaltsverzeichnis einfügen
\section*{Internetquellen}
\printbibliography[type=online,heading=none,title={Internetquellen}]%Literaturverzeichnis der Onlinequellen drucken.
%--------FÜR HAUSARBEIT BEI KRAUSE----------------------
% \addchap{Literaturverzeichnis}
% \printbibliography[nottype=online, nottype=article,heading=none]%Literaturverzeichnis der NICHT-Artikel/Onlinequellen drucken.
% \addchap{Quellenverzeichnis}
% \printbibliography[filter=articleonline,heading=none]%Literaturverzeichnis der /ArtikelOnlinequellen drucken.
%-------------------------------------------------------
% !TeX root = ../../main.tex

%-----------------------------------
% Ehrenwörtliche Erklärung
%-----------------------------------
\chapter*{Ehrenwörtliche Erklärung}
\pagenumbering{gobble} % Keine Seitenzahlen mehr
\renewcommand*\chapterpagestyle{leer}
Hiermit versichere ich, dass die vorliegende Arbeit von mir selbstständig und ohne unerlaubte Hilfe angefertigt worden ist, insbesondere dass ich alle Stellen, die wörtlich oder annähernd wörtlich aus Veröffentlichungen entnommen sind, durch Zitate als solche gekennzeichnet habe. Ich versichere auch, dass die von mir eingereichte schriftliche Version mit der digitalen Version übereinstimmt. Weiterhin erkläre ich, dass die Arbeit in gleicher oder ähnlicher Form noch keiner Prüfungsbehörde/Prüfungsstelle vorgelegen hat. Ich erkläre mich damit einverstanden, dass die Arbeit der Öffentlichkeit zugänglich gemacht wird. Ich erkläre mich damit einverstanden, dass die Digitalversion dieser Arbeit zwecks Plagiatsprüfung auf die Server externer Anbieter hoch geladen werden darf. Die Plagiatsprüfung stellt keine Zurverfügungstellung für die Öffentlichkeit dar.

\par\medskip
\par\medskip

\vspace{5cm}

\begin{table}[H]
	\centering
	\begin{tabular*}{\textwidth}{c @{\extracolsep{\fill}} ccccc}
		\myOrt, \today
		&
		% Hinterlege deine eingescannte Unterschrift im Verzeichnis /abbildungen und nenne sie unterschrift.png
		% Bilder mit transparentem Hintergrund können teils zu Problemen führen
		\includegraphics[width=0.35\textwidth]{Abbildungen/unterschrift.png}\vspace*{-0.35cm}
		\\
		\rule[0.5ex]{12em}{0.55pt} & \rule[0.5ex]{12em}{0.55pt} \\
		(Ort, Datum) & (Eigenhändige Unterschrift)
		\\
	\end{tabular*} \\
\end{table}
\end{document}%Ende des Dokumentes
